%=============================================================
% CHAPTER 1 — SOFTWARE PROJECT MANAGEMENT PLAN (SPMP)
%=============================================================

\chapter{Software Project Management Plan (SPMP)}

%-------------------------------------------------------------
\section{Introduction}
%-------------------------------------------------------------

Pakistan is home to over 230 million people, and for most of them Urdu or Roman Urdu is
the language they actually think in. Yet if someone in Lahore or Karachi wants to look up
drug interactions or understand a diagnosis, every halfway-decent AI tool on the internet
responds in one language only: English. Symptom checkers, clinical chatbots, question-answering
systems --- built for English, returning English, assuming English is the default.

That is not a small gap. For a population this size, it is a structural barrier to basic
health information. General-purpose AI assistants like ChatGPT can technically respond in
Urdu, but they pull their medical knowledge straight from training weights --- no verified
source, no updated corpus, no way to trace where a claim came from. In a medical setting,
a confident but wrong answer is not just unhelpful. It is dangerous.

What the problem needs is a system that combines the language flexibility of modern LLMs
with retrieval from a trusted, curated medical knowledge base. Figure~\ref{fig:system_overview}
shows the high-level design of MediLingo AI, which connects a React/Vite frontend to a
FastAPI backend running a full RAG pipeline --- PDF ingestion, semantic indexing, and
multilingual response generation via a hosted large language model.

\begin{figure}[H]
  \centering
  \includegraphics[width=0.90\textwidth]{system_overview}
  \caption{MediLingo AI --- High-Level System Overview}
  \label{fig:system_overview}
\end{figure}

%-------------------------------------------------------------
\subsection{Existing Medical Information Systems}
%-------------------------------------------------------------

\subsubsection{Existing Systems With AI Capabilities}

A number of commercial platforms already use AI for medical information delivery. Each one
has the same blind spot: they are built around English, and there is no real pathway for
someone who wants to type a query in Urdu.

\begin{itemize}

  \item \textbf{Ada Health \cite{adahealth2024}:} Ada is a conversational symptom checker
        that walks users through a structured set of questions and returns a ranked
        differential of probable conditions. Its clinical reasoning engine is proprietary
        and English-only. It does not accept free-text queries in Urdu and provides no
        retrieval mechanism against a user-supplied knowledge base.

        \begin{figure}[H]
          \centering
          \figplaceholder[4cm]{Screenshot of Ada Health symptom checker interface}
          \caption{Ada Health --- AI-driven symptom checker (English only)}
          \label{fig:ada}
        \end{figure}

  \item \textbf{Babylon Health:} Babylon's AI triage service operates on a closed clinical
        model trained on English-language medical data. It has never been deployed in
        Pakistan and does not support Urdu as either an input or output language.

        \begin{figure}[H]
          \centering
          \figplaceholder[4cm]{Screenshot of Babylon Health AI consultation interface}
          \caption{Babylon Health --- AI triage and health consultation}
          \label{fig:babylon}
        \end{figure}

  \item \textbf{Med-PaLM 2 (Google) \cite{singhal2023medpalm}:} Med-PaLM 2 is a large
        language model fine-tuned on medical benchmarks. It scored at expert level on
        medical licensing exams, which is impressive. But it is a research system, not
        publicly available, and it operates in English only --- so for a Urdu-speaking
        user, its benchmark scores are essentially irrelevant.

  \item \textbf{WebMD Symptom Checker:} WebMD uses a keyword-driven symptom tool that
        matches inputs to conditions in a static database. No natural-language generation,
        no retrieval pipeline, no Urdu support. It is more of a lookup tool than an
        intelligent assistant.

\end{itemize}

\subsubsection{Existing Systems Without Multilingual AI}

On the Pakistani side, the platforms that actually serve this population provide health
information without any real AI query capability:

\begin{itemize}

  \item \textbf{Sehat.com.pk and Marham.pk:} These are the most widely used Pakistani
        health portals. They list doctors, allow appointment booking and publish basic
        health articles --- some of it in Urdu. But neither offers an interactive medical
        query interface, and there is no generative response system or retrieval from a
        knowledge base. Useful for finding a doctor; not useful for asking a clinical
        question.

  \item \textbf{General-purpose LLMs (e.g., ChatGPT):} Models like ChatGPT do understand
        Urdu and Roman Urdu reasonably well. The problem is that their medical knowledge
        comes from training data, not a verified clinical corpus. It cannot be updated, and
        there is no way to trace a claim back to a source document. That is a meaningful
        risk in a medical context --- the answer might be fluent and confident while being
        clinically wrong or outdated.

  \item \textbf{Government and hospital websites:} Static health content published by
        Pakistani hospitals and the Ministry of National Health Services is in English.
        None of these platforms offer any interactive query interface at all.

\end{itemize}

%-------------------------------------------------------------
\subsection{Project Overview}
%-------------------------------------------------------------

MediLingo AI is a web-based medical chatbot that takes queries in English, Urdu and Roman
Urdu and returns responses drawn from an indexed collection of medical PDF documents. The
key distinction from a general-purpose chatbot is that responses are not generated from
model memory --- they are grounded in passages retrieved from the vector store at query
time. The LLM is there to synthesise and explain, not to recall.

The backend runs on FastAPI \cite{fastapi2018} and handles four things: user authentication
and session storage via MongoDB \cite{mongodb2024} and JWT tokens; PDF ingestion and
semantic chunking via \texttt{pdfplumber}; dense retrieval via ChromaDB \cite{chromadb2023}
and \texttt{sentence-transformers} \cite{sentencetransformers2019}; and response generation
through the Groq-hosted LLaMA~3.3-70b model \cite{groq2024}. The frontend is React~19 with
TypeScript and Vite \cite{react2024}, styled with TailwindCSS, and uses the browser's Web
Speech API for voice interaction.

The key features of the system are:

\begin{itemize}
  \item User registration, login and password reset via email OTP
  \item Multilingual query input in English, Urdu and Roman Urdu
  \item RAG-grounded responses generated from ingested medical PDFs
  \item Full chat session management: create, continue, rename and delete sessions
  \item Persistent message history stored in MongoDB
  \item Voice input (speech-to-text) and voice output (text-to-speech)
  \item Admin panel for user management and triggering PDF ingestion
  \item JWT dual-token authentication with HTTP-only refresh cookies
\end{itemize}

%-------------------------------------------------------------
\subsection{Project Deliverables}
%-------------------------------------------------------------

\begin{itemize}
  \item Software Project Management Plan (SPMP)
  \item Software Requirements Specification (SRS)
  \item Software Design Description (SDD)
  \item Software Test Documentation (STD)
  \item FastAPI backend with integrated RAG engine and ingestion pipeline
  \item React / TypeScript / Vite frontend with multilingual chat interface
  \item ChromaDB vector store populated from medical PDF documents
  \item MongoDB schema for users, sessions and messages
  \item Docker and Docker Compose configuration for single-command deployment
  \item GitHub repository with complete version-controlled source code
  \item Final working demonstration of the cross-lingual RAG flow
\end{itemize}

%=============================================================
\section{Project Organization}
%=============================================================

%-------------------------------------------------------------
\subsection{Software Process Model}
%-------------------------------------------------------------

The \textbf{Iterative Incremental Model} was chosen for MediLingo AI. The reasoning was
straightforward once the system's structure was clear:

\begin{itemize}

  \item The system breaks down naturally into independent, shippable increments.
        Authentication and session storage can be built and fully verified before a single
        line of RAG code is written. The English-only RAG pipeline can be validated before
        Urdu support is layered on. Voice features can be added last without touching
        anything below them. Each increment leaves the system in a working state.

  \item A RAG pipeline requires empirical tuning --- chunk size, overlap length, top-k
        retrieval count, similarity threshold. These values cannot be derived analytically;
        they have to be evaluated against real queries. An iterative model makes this
        feedback loop part of the process rather than something tacked on at the end.

  \item The language-prefix stripping mechanism was added mid-project, after it became
        clear that embedding the full prefixed query was degrading retrieval quality. An
        iterative approach meant this design change could be isolated to one increment,
        tested and confirmed without breaking anything already shipped.

\end{itemize}

Figure~\ref{fig:process_model} shows the four increments and how each builds on the last.

\begin{figure}[H]
  \centering
  \includegraphics[width=0.75\textwidth]{process_model}
  \caption{Iterative Incremental Software Process Model applied to MediLingo AI}
  \label{fig:process_model}
\end{figure}

%-------------------------------------------------------------
\subsection{Roles and Responsibilities}
%-------------------------------------------------------------

\textbf{Developer (Student):}
\begin{itemize}
  \item Design and implement all FastAPI routers covering authentication
        (\texttt{/auth/*}), session management (\texttt{/sessions/*}) and RAG query
        processing (\texttt{/query}, \texttt{/retrieve}, \texttt{/pipeline/*})
  \item Build the RAG ingestion pipeline: PDF extraction with \texttt{pdfplumber},
        semantic chunking, embedding generation and ChromaDB indexing
  \item Integrate the Groq API and implement the language-prefix stripping mechanism so
        that semantic search operates on clean English text regardless of the query language
  \item Develop the React frontend including the chat page, session sidebar, login,
        register, profile, admin, forgot-password and reset-password pages
  \item Implement voice input and output using the Web Speech API through custom hooks
  \item Set up JWT dual-token authentication with silent refresh on 401 responses
  \item Configure Docker and Docker Compose for the full stack
  \item Write all project documentation and prepare diagrams
\end{itemize}

\textbf{Supervisor:}
\begin{itemize}
  \item Provide architectural guidance on the RAG design and multilingual handling strategy
  \item Review documentation milestones and provide written feedback
  \item Clarify requirements and resolve design ambiguities during development
\end{itemize}

%-------------------------------------------------------------
\subsection{Tools and Techniques}
%-------------------------------------------------------------

\textbf{Tools Used:}

\begin{itemize}
  \item \textbf{Python 3.11 and FastAPI \cite{fastapi2018}} --- asynchronous REST API
        framework with built-in OpenAPI documentation at \texttt{/docs}
  \item \textbf{React 19, TypeScript and Vite \cite{react2024}} --- component-based
        frontend with fast hot-module replacement during development
  \item \textbf{TailwindCSS} --- utility-first CSS framework for all UI styling
  \item \textbf{Zustand} --- minimal state management for auth tokens and session list
  \item \textbf{MongoDB and Motor \cite{mongodb2024}} --- async NoSQL store for
        \texttt{UserDoc}, \texttt{SessionDoc} and \texttt{MessageDoc} records
  \item \textbf{ChromaDB \cite{chromadb2023}} --- embedded vector database; persists to
        \texttt{backend/vectorstore/} and supports cosine similarity search
  \item \textbf{sentence-transformers \texttt{all-MiniLM-L6-v2}
        \cite{sentencetransformers2019}} --- 384-dimensional dense embedding model loaded
        once at startup via the RAG service singleton
  \item \textbf{Groq API (LLaMA~3.3-70b) \cite{groq2024}} --- cloud-hosted inference
        endpoint for final response generation
  \item \textbf{pdfplumber / pypdf} --- PDF text extraction for the ingestion pipeline
  \item \textbf{Docker and Docker Compose} --- containerised deployment of backend,
        frontend and MongoDB as a single \texttt{docker-compose up} command
  \item \textbf{Git and GitHub} --- version control; all increments developed on the same
        main branch with descriptive commits
  \item \textbf{PlantUML} --- UML diagram generation for all SRS and SDD sequence diagrams,
        domain model, class diagram, use case diagram and Gantt chart
  \item \textbf{VS Code + MiKTeX + LaTeX Workshop} --- documentation typesetting for all
        four documents compiled locally with pdflatex
\end{itemize}

\textbf{Techniques Used:}

\begin{itemize}
  \item Retrieval-Augmented Generation (RAG) \cite{lewis2020rag} --- retrieval from a
        curated domain corpus before language model generation
  \item Semantic chunking at 1000 characters with 200-character overlap to preserve
        sentence context across chunk boundaries
  \item Cosine similarity search over 384-dimensional embeddings for passage retrieval
  \item Language-prefix system: user language instruction prepended to query for the LLM
        but stripped before embedding to keep retrieval language-neutral
  \item JWT dual-token authentication: 15-minute access tokens (in-memory only) plus
        7-day refresh tokens (HTTP-only cookie) with automatic silent refresh on 401
  \item RAG singleton lifecycle managed by \texttt{rag\_service.py} with deferred
        import of heavy ML libraries (\texttt{sentence-transformers}, \texttt{chromadb})
        to avoid slowing FastAPI startup
  \item UUID4 string IDs for all MongoDB documents to decouple application identity from
        the MongoDB \texttt{\_id} field
\end{itemize}

%=============================================================
\section{Project Management Plan}
%=============================================================

%-------------------------------------------------------------
\subsection{Tasks}
%-------------------------------------------------------------

The project is organised into six phases executed in iterative sequence:

\begin{itemize}
  \item Requirements Analysis Phase
  \item System Design Phase
  \item Backend Development Phase
  \item Frontend Development Phase
  \item RAG Integration and Multilingual Testing Phase
  \item Testing and Validation Phase
\end{itemize}

%-------------------------------------------------------------
\subsection{Description}
%-------------------------------------------------------------

The \textbf{Requirements Analysis Phase} was really about pinning down what the system
needed to do before any code was written. This meant listing out the functional
requirements --- registration, multilingual medical query, session management, voice I/O,
admin ingestion trigger --- and the non-functional ones around latency, token security and
browser compatibility for voice features.

During the \textbf{System Design Phase}, the three-tier architecture was finalised
alongside the MongoDB document schema, the ChromaDB collection design and the RAG pipeline
dataflow. The language-prefix stripping contract between frontend and backend was one of
the more important design decisions made here. UML artefacts --- use-case diagrams, system
sequence diagrams, domain model and class diagram --- came out of this phase.

The \textbf{Backend Development Phase} is where the actual implementation began. All
FastAPI routers were built, along with \texttt{RAGQueryEngine}, the ingestion pipeline
class \texttt{RAGIngestionPipeline}, the Motor-based database layer and the JWT auth
service. The Groq API integration and the async lifespan startup sequence were also
completed here.

The \textbf{Frontend Development Phase} covered all the React pages and custom hooks:
landing page, chat interface with session sidebar, login and registration flows, profile
page, forgot-password and reset-password pages, admin panel, and the
\texttt{useVoiceInput} / \texttt{useVoiceOutput} hooks. Vite's proxy was configured so
API routes forwarded to \texttt{localhost:8000} during development --- this was important
for getting cookies to work correctly without fighting CORS.

The \textbf{RAG Integration and Multilingual Testing Phase} connected the live frontend
to the backend and ran queries in all three languages. This was the phase where the
language-prefix mechanism was validated --- the test was simple: switching language should
not change which documents get retrieved, only how the response is worded. It passed.

The \textbf{Testing and Validation Phase} covered functional testing of every use case:
JWT token lifecycle, RAG retrieval quality, voice API in Chrome and Edge, and admin
workflows. Issues found here were fixed before final submission.

%-------------------------------------------------------------
\subsection{Project Plan}
%-------------------------------------------------------------

Figure~\ref{fig:gantt} shows the project plan as a Gantt chart. The backend and frontend
development phases overlap intentionally --- waiting for the backend to be completely
finished before touching the frontend would have stretched the timeline by weeks. The
RAG integration phase that follows both is where the two halves are connected and
tested end-to-end.

\begin{figure}[H]
  \centering
  \includegraphics[width=\textwidth]{gantt}
  \caption{Project Management Plan --- Gantt Chart}
  \label{fig:gantt}
\end{figure}

%-------------------------------------------------------------
\subsection{Project Timeline}
%-------------------------------------------------------------

Figure~\ref{fig:timeline} marks the key milestones from requirements analysis to final
submission. The documentation-heavy first half and implementation-heavy second half are
not a coincidence --- writing down what the system needs to do before writing the code
for it saves significant rework later. That principle drove the schedule.

\begin{figure}[H]
  \centering
  \includegraphics[width=0.45\textwidth]{timeline}
  \caption{MediLingo AI --- Project Timeline}
  \label{fig:timeline}
\end{figure}
